\documentclass[letterpaper]{article}
\usepackage[margin=1in]{geometry}
\usepackage{times}
\usepackage{babel}
\usepackage[T1]{fontenc}
\usepackage[utf8]{inputenc}
\usepackage{amsmath,amssymb,amsthm}
\usepackage{graphicx}
\usepackage{booktabs}
\usepackage{array}
\usepackage{xcolor}
\usepackage[colorlinks=true,linkcolor=blue,citecolor=green!60!black,urlcolor=blue]{hyperref}
\usepackage{natbib}
\usepackage{microtype}
\usepackage{enumitem}
\usepackage{float}
\usepackage[ruled,vlined,linesnumbered]{algorithm2e}
\usepackage{caption}

\newtheorem{theoreme}{Th\'eor\`eme}
\newtheorem{corollaire}{Corollaire}[theoreme]
\newtheorem{definition}{D\'efinition}
\newtheorem{proposition}{Proposition}
\newtheorem{remarque}{Remarque}

\title{%
\textbf{Inf\'erence d'\'Etat de R\`egles (RSI) :}\\
\textbf{Un Cadre Bay\'esien pour la Surveillance de la Conformit\'e}\\
\textbf{dans les Domaines R\'egis par des R\`egles}\\[0.4em]
\large\textit{Application aux Syst\`emes Fiscaux d'Afrique Francophone}
}
\author{
    \textbf{Abdou-Raouf Atarmla}$^{1,2}$\thanks{Travail ind\'ependant. Correspondance : \texttt{achilleatarmla@gmail.com}} \\
    \small $^{1}$Institut National des Postes et T\'el\'ecommunications (INPT), Rabat, Maroc \\
    \small $^{2}$Togo DataLab, Minist\`ere de l'\'Economie Num\'erique, Lom\'e, Togo \\
    \small \textbf{Emails :} \texttt{achilleatarmla@gmail.com} $|$ \texttt{abdou-raouf.atarmla@datalab.gouv.tg} \\
    \small \texttt{atarmla.abdouraouf@ine.inpt.ac.ma}
}
\date{2026}

\begin{document}
\maketitle
\thispagestyle{empty}

%% ─────────────────────────────────────────────────────────────────
\begin{abstract}
\noindent
Les cadres d'apprentissage automatique existants pour la surveillance
de la conformit\'e---r\'eseaux logiques de Markov, logique floue
probabiliste, mod\`eles supervis\'es---partagent un paradigme fondamental~:
ils traitent les donn\'ees observ\'ees comme une v\'erit\'e de terrain
et tentent d'approximer les r\`egles \`a partir de celles-ci.
Cette hypoth\`ese s'effondre dans les \emph{domaines r\'egis par des
r\`egles}, tels que la fiscalit\'e ou la conformit\'e r\'eglementaire,
o\`u les r\`egles autorit\'es sont \emph{connues a priori} et o\`u le
v\'eritable d\'efi consiste \`a inf\'erer l'\emph{\'etat latent} d'activation
des r\`egles, de conformit\'e et de d\'erive param\'etrique \`a partir
d'observations partielles et bruit\'ees.

Nous proposons \textbf{l'Inf\'erence d'\'Etat de R\`egles (RSI)}, un
cadre bay\'esien qui inverse ce paradigme en codant les r\`egles
r\'eglementaires comme des a priori structur\'es et en formulant la
surveillance de la conformit\'e comme une inf\'erence post\'erieure sur un
espace d'\'etats latents $\mathbf{S} = \{(a_i,\, c_i,\, \delta_i)\}_{i=1}^n$,
o\`u $a_i$ capture l'activation de la r\`egle, $c_i$ mod\'elise le taux
de conformit\'e, et $\delta_i$ quantifie la d\'erive param\'etrique.
Nous prouvons trois garanties th\'eoriques~:
\textbf{(T1)} RSI absorbe les changements r\'eglementaires en temps $O(1)$
via une correction de ratio d'a priori, ind\'ependamment de la taille
du jeu de donn\'ees~;
\textbf{(T2)} le posterior est consistent au sens de Bernstein-von Mises,
convergeant vers le vrai \'etat des r\`egles \`a mesure que les
observations s'accumulent~;
\textbf{(T3)} l'inf\'erence variationnelle champ-moyen maximise
monotonement la borne inf\'erieure de l'\'evidence (ELBO).

Nous instancions RSI sur le syst\`eme fiscal togolais et introduisons
\textbf{RSI-Togo-Fiscal-Synthetic v1.0}, un jeu de donn\'ees de r\'ef\'erence
de 2~000 entreprises synth\'etiques ancr\'e dans les r\`egles r\'eelles de
l'OTR (2022--2025). Sans aucune donn\'ee d'entra\^inement \'etiquet\'ee,
RSI atteint F1~$=$~0,519 et AUC~$=$~0,599, tout en absorbant les changements
r\'eglementaires en 0,2~ms contre 407~ms pour le r\'e-entra\^inement
complet---une acc\'el\'eration de $2~000\times$.

\medskip
\noindent\textbf{Mots-cl\'es :} inf\'erence bay\'esienne, domaines r\'egis
par des r\`egles, surveillance de la conformit\'e, inf\'erence variationnelle,
raisonnement z\'ero-shot, conformit\'e fiscale, Afrique
\end{abstract}

\newpage

%% ─────────────────────────────────────────────────────────────────
\section{Introduction}
\label{sec:intro}

L'apprentissage automatique a r\'ealis\'e des progr\`es remarquables
en reconnaissance de formes, pr\'evision et classification. Pourtant,
dans une large classe de domaines \`a forts enjeux---administration
fiscale, conformit\'e m\'edicale, r\'eglementation juridique, audit
financier---les praticiens se heurtent \`a une difficult\'e structurelle
r\'ecurrente~: \emph{la connaissance est d\'ej\`a structur\'ee}. Les
r\`egles ne sont pas inconnues. Elles sont \'ecrites dans la loi, codifi\'ees
dans la r\'eglementation et appliqu\'ees par des institutions.

Dans ces \textbf{domaines r\'egis par des r\`egles}, le d\'efi n'est pas
d'apprendre les r\`egles \`a partir des donn\'ees. Le d\'efi est d'inf\'erer
\emph{si les entit\'es les respectent}, sachant que les observations sont
partielles, strat\'egiquement manipul\'ees et fr\'equemment manquantes.
Un inspecteur fiscal n'a pas besoin d'un r\'eseau de neurones pour
d\'eterminer qu'une entreprise avec 80 millions de FCFA de chiffre
d'affaires est assujettie \`a la TVA---la loi le stipule. Ce dont il a
besoin, c'est d'une m\'ethode rigoureuse pour \'evaluer, \`a partir d'un
ensemble incomplet de d\'eclarations et de signaux comportementaux,
l'\'etat r\'eel de conformit\'e de cette entreprise.

\medskip
\noindent\textbf{Limites des approches existantes.}
Les classificateurs supervis\'es (XGBoost, r\'eseaux de neurones)
n\'ecessitent des exemples \'etiquet\'es de non-conformit\'e, qui sont
rares, juridiquement sensibles et souvent indisponibles en pratique.
Les r\'eseaux logiques de Markov (MLN)~\citep{richardson2006markov} et
la logique floue probabiliste (PSL)~\citep{bach2017hinge} traitent les
poids des r\`egles comme des quantit\'es \`a \emph{apprendre} \`a partir
des donn\'ees, inversant la direction \'epist\'emique. Les m\'ethodes
neurosymboliques~\citep{garcez2022neural} combinent perception neuronale
et raisonnement symbolique, mais leur composante symbolique demeure une
approximation apprise. Aucun de ces cadres ne fournit~: (i)~une
\'evaluation de conformit\'e z\'ero-shot \`a partir de r\`egles autorit\'ees,
(ii)~une adaptation $O(1)$ aux changements r\'eglementaires, ni
(iii)~une quantification native de l'incertitude interpr\'etable par
des auditeurs non techniques.

\medskip
\noindent\textbf{Nos contributions.}
Nous proposons \textbf{RSI (Rule-State Inference)}, un cadre bay\'esien
qui formalise le paradigme inverse. Nos contributions sont~:
\begin{enumerate}[leftmargin=2em, label=(\arabic*)]
  \item \textbf{Le cadre RSI}~: une formalisation rigoureuse de la
        surveillance de la conformit\'e comme inf\'erence post\'erieure
        bay\'esienne sur un espace d'\'etats latents, avec les r\`egles
        comme a priori structur\'es (Section~\ref{sec:framework}).
  \item \textbf{Trois garanties th\'eoriques}~: adaptabilit\'e $O(1)$
        aux changements r\'eglementaires (T1), consistance post\'erieure
        de Bernstein-von Mises (T2), et convergence monotone de l'ELBO
        (T3) (Section~\ref{sec:theory}).
  \item \textbf{RSI-Togo-Fiscal-Synthetic v1.0}~: un jeu de donn\'ees
        de r\'ef\'erence publiquement disponible, ancr\'e dans la loi
        fiscale togolaise, avec un \'ev\'enement de changement
        r\'eglementaire document\'e (Section~\ref{sec:experiments}).
  \item \textbf{Validation empirique} \`a travers cinq exp\'eriences
        d\'emontrant des performances z\'ero-shot, une acc\'el\'eration
        $2~000\times$ de mise \`a jour, et une robustesse \`a 50\% de
        donn\'ees manquantes (Section~\ref{sec:experiments}).
\end{enumerate}

\medskip
\noindent\textbf{Pourquoi l'Afrique francophone.}
Nous ancrons notre cadre dans le syst\`eme fiscal togolais pour trois
raisons qui g\'en\'eralisent largement. Premi\`erement, la fr\'equence
des changements r\'eglementaires est \'elev\'ee (plusieurs amendements
par ann\'ee fiscale), rendant l'adaptabilit\'e non pas un luxe mais
une n\'ecessit\'e. Deuxi\`emement, la qualit\'e des donn\'ees est
structurellement faible~: les revenus d\'eclar\'es sont syst\'ematiquement
sous-d\'eclar\'es (ratio moyen de 0,70 dans notre benchmark), et 18--20\%
des variables fiscales cl\'es sont manquantes. Troisi\`emement, les
solutions ML d\'evelopp\'ees pour les contextes \`a revenus \'elev\'es
supposent une richesse de donn\'ees qui n'existe tout simplement pas dans
ces environnements. RSI est con\c{c}u \emph{\`a partir} de ces contraintes,
non adapt\'e \emph{pour} elles.

%% ─────────────────────────────────────────────────────────────────
\section{Travaux connexes}
\label{sec:related}

\paragraph{Logique probabiliste.}
Les MLN~\citep{richardson2006markov} combinent la logique du premier
ordre avec les champs al\'eatoires de Markov via des formules pond\'er\'ees,
avec des poids appris \`a partir des donn\'ees.
La PSL~\citep{bach2017hinge} utilise des valeurs de v\'erit\'e continues
et r\'ealise une inf\'erence MAP tractable via l'optimisation convexe.
Les deux traitent les poids des r\`egles comme des \emph{sorties} de
l'apprentissage. RSI traite les r\`egles comme des \emph{entr\'ees}~:
leur statut \'epist\'emique est a priori, non approxim\'e.
Dans les MLN, une formule de poids z\'ero est effectivement absente~;
dans RSI, une r\`egle avec a priori $\pi_i = 0,99$ est quasi-certaine
et prime sur les preuves observationnelles faibles.

\paragraph{IA neurosymbolique.}
Garcez et Lamb~\citeyearpar{garcez2022neural} d\'ecrivent des architectures
allant des mod\`eles en pipeline aux syst\`emes enti\`erement
diff\'erentiables (Neural Theorem Provers~\citep{rocktaschel2017end},
LTN~\citep{badreddine2022logic}). La direction canonique est
l'\emph{apprentissage pour le raisonnement}~: la structure symbolique
guide l'apprentissage neuronal. RSI occupe la position orthogonale~:
\emph{raisonnement depuis les r\`egles}, o\`u la couche symbolique est
fix\'ee et autorit\'ee, non apprise.

\paragraph{Apprentissage interpr\'etable de r\`egles.}
Les Listes de R\`egles Bay\'esiennes~\citep{letham2015interpretable}
appliquent la s\'election de mod\`eles bay\'esienne \`a des ensembles
de r\`egles---encore une fois, les r\`egles sont la sortie.
Notre travail produit non pas des r\`egles mais des \emph{distributions
post\'erieures sur les \'etats de conformit\'e}---une cible d'inf\'erence
qualitativement diff\'erente.

\paragraph{IA et conformit\'e fiscale.}
Le ML a \'et\'e appliqu\'e \`a l'estimation de l'\'ecart fiscal
~\citep{gomes2022machine}, \`a la s\'election d'audits et \`a la
d\'etection de fraudes TVA, universellement sous des paradigmes
supervis\'es. Aucun travail ant\'erieur ne formalise la surveillance
de la conformit\'e comme une inf\'erence d'\'etat bay\'esienne.

\paragraph{Donn\'ees manquantes et apprentissage en faible ressource.}
Fr\'enay et Verleysen~\citeyearpar{frenay2014classification} recensent
l'apprentissage en pr\'esence de bruit de label. RSI adresse un d\'efi
connexe mais distinct~: l'absence \emph{structurelle} de donn\'ees dans
des environnements \`a faible int\'egrit\'e. Notre vraisemblance g\`ere
les valeurs manquantes nativement en assignant une probabilit\'e
unitaire, pr\'eservant l'a priori sans injection de biais.

%% ─────────────────────────────────────────────────────────────────
\section{Le Cadre RSI}
\label{sec:framework}

\subsection{Formulation du probl\`eme}

Soit $\mathcal{R} = \{r_1, \ldots, r_n\}$ un ensemble de r\`egles
r\'eglementaires connues. Chaque r\`egle $r_i : \mathcal{X} \times
\Theta_i \to \{0,1\}$ fait correspondre les attributs d'une entit\'e
$x \in \mathcal{X}$ et les param\`etres de la r\`egle $\Theta_i$ \`a
un jugement binaire d'applicabilit\'e. Soit
$\mathbf{D} = \{d_1, \ldots, d_m\}$ un ensemble d'observations bruit\'ees
du comportement des entit\'es.

\begin{definition}[Domaine r\'egi par des r\`egles]
Un domaine $(\mathcal{R}, \mathcal{X}, \mathbf{D})$ est \emph{r\'egi
par des r\`egles} si~: (i)~$\mathcal{R}$ est autorit\'e et connue a
priori~; (ii)~$\mathbf{D}$ est partiel, potentiellement manquant, et
peut \^etre strat\'egiquement distorc\'e~; (iii)~la t\^ache d'inf\'erence
principale est d'\'evaluer la conformit\'e des entit\'es avec
$\mathcal{R}$, non de l'apprendre.
\end{definition}

\subsection{L'espace d'\'etats latents}

\begin{definition}[\'Etat de r\`egle]
L'\'etat latent de la r\`egle $r_i$ est~:
\[
  s_i = (a_i,\; c_i,\; \delta_i), \quad
  \mathbf{S} = (s_1, \ldots, s_n) \in \mathcal{S}
\]
o\`u $a_i \in \{0,1\}$ indique si $r_i$ est en vigueur, $c_i \in
[0,1]$ est le taux de conformit\'e, et $\delta_i \in \mathbb{R}$
capture la d\'erive param\'etrique.
\end{definition}

\noindent\textbf{Intuition fondamentale.} Cet objet n'a pas
d'\'equivalent dans les cadres ML existants. Les MLN ont des poids de
r\`egles~; RSI a des \'etats de r\`egles. La diff\'erence est
\'epist\'emique~: les poids sont inf\'er\'es \`a partir des donn\'ees
(donn\'ees~$\to$~r\`egle)~; les \'etats sont inf\'er\'es \'etant donn\'e
les r\`egles (r\`egles + donn\'ees~$\to$~\'etat).

\subsection{L'a priori $P(\mathbf{S})$ : les r\`egles comme a priori structur\'es}

Sous la factorisation champ-moyen~:
\[
  P(\mathbf{S}) = \prod_{i=1}^n P(a_i)\, P(c_i \mid a_i)\, P(\delta_i \mid a_i)
\]
\begin{align}
  a_i &\sim \operatorname{Bernoulli}(\pi_i) \\
  c_i \mid a_i{=}1 &\sim \operatorname{Beta}(\alpha_i, \beta_i) \\
  \delta_i \mid a_i{=}1 &\sim \mathcal{N}(0,\, \sigma_i^2)
\end{align}
Les hyperparams $(\pi_i, \alpha_i, \beta_i, \sigma_i)$ encodent la
connaissance institutionnelle~: taux de conformit\'e historiques,
a priori sur l'activation des r\`egles, stabilit\'e param\'etrique
attendue.

\subsection{La vraisemblance $P(\mathbf{D} \mid \mathbf{S})$}

Les observations sont conditionnellement ind\'ependantes~:
\[
  P(\mathbf{D} \mid \mathbf{S}) = \prod_{j=1}^m P(d_j \mid \mathbf{S})
\]
Pour chaque observation $d_j$~:
\[
  P(d_j \mid \mathbf{S}) = \sum_{i:\, r_i \text{ applicable}}
  \Bigl[a_i \cdot \mathcal{L}(d_j \mid r_i, c_i, \delta_i)
  + (1 - a_i) \cdot \varepsilon\Bigr]
\]
\textbf{Gestion des donn\'ees manquantes~:} quand $d_j$ est absent,
$P(d_j \mid \mathbf{S}) \equiv 1$, pr\'eservant l'a priori.

Nous utilisons des fonctions de vraisemblance \emph{segment\'ees} par
r\'egime fiscal~: (i)~\textbf{TPU} (secteur informel, CA~$<$~30M FCFA)~:
signaux~: d\'elai de paiement, compte bancaire, ratio de
sous-d\'eclaration~; (ii)~\textbf{TVA} (CA~$\geq$~seuil)~: TVA
d\'eclar\'ee vs th\'eorique, mod\`ele bruit log-gaussien~;
(iii)~\textbf{IS} (CA~$\geq$~100M FCFA)~: IS d\'eclar\'e vs IS
th\'eorique donn\'e le b\'en\'efice d\'eclar\'e.

\subsection{Inf\'erence post\'erieure}

\[
  \boxed{P(\mathbf{S} \mid \mathbf{D}) \;\propto\;
    P(\mathbf{D} \mid \mathbf{S}) \cdot P(\mathbf{S})}
\]

Le posterior supporte trois requ\^etes interpr\'etables~:
\begin{itemize}[leftmargin=1.5em]
  \item $P(a_i{=}1 \mid \mathbf{D})$~: probabilit\'e que $r_i$ soit active
  \item $\mathbb{E}[c_i \mid \mathbf{D}] \pm \sqrt{\mathrm{Var}[c_i \mid \mathbf{D}]}$~:
        conformit\'e attendue avec incertitude
  \item $\mathbb{E}[\delta_i \mid \mathbf{D}]$~: d\'erive param\'etrique estim\'ee
\end{itemize}
Toutes les quantit\'es sont directement interpr\'etables par des auditeurs
non techniques, sans outil d'explication post-hoc.

\subsection{Inf\'erence variationnelle champ-moyen}

L'inf\'erence exacte est intractable pour $n$ grand. Nous approximons
le posterior en minimisant $\mathrm{KL}[Q_\phi(\mathbf{S}) \|
P(\mathbf{S} \mid \mathbf{D})]$, \'equivalent \`a maximiser~:
\[
  \mathrm{ELBO}(\phi) = \mathbb{E}_{Q_\phi}[\log P(\mathbf{D} \mid \mathbf{S})]
  - \mathrm{KL}[Q_\phi(\mathbf{S}) \| P(\mathbf{S})]
\]
Toutes les mises \`a jour sont \emph{analytiques} (familles conjugu\'ees)~:
Beta-Binomiale pour $c_i$, Bernoulli pour $a_i$, Gaussienne pour $\delta_i$.

\begin{algorithm}[H]
\SetAlgoLined
\KwIn{Observations $\mathbf{D}$, R\`egles $\mathcal{R}$, A priori
      $\{\pi_i, \alpha_i, \beta_i, \sigma_i\}$}
\KwOut{Posterior $\{P(a_i|\mathbf{D}),\; \mathbb{E}[c_i|\mathbf{D}],\;
       \mathbb{E}[\delta_i|\mathbf{D}]\}$}
Initialiser $Q_\phi \leftarrow P(\mathbf{S})$\;
\Repeter{$|\mathrm{ELBO}_t - \mathrm{ELBO}_{t-1}| < \varepsilon$}{
  \PourCh{r\`egle $r_i \in \mathcal{R}$}{
    Calculer les signaux de conformit\'e depuis $\mathbf{D}$\;
    Mettre \`a jour $Q(a_i) \leftarrow \operatorname{Bernoulli}(\rho_i)$\;
    Mettre \`a jour $Q(c_i) \leftarrow \operatorname{Beta}(\alpha_i + n_{\mathrm{ok}},\;
                              \beta_i + n_{\mathrm{\'echec}})$\;
    Mettre \`a jour $Q(\delta_i) \leftarrow \mathcal{N}(\mu_{\mathrm{post}},
                               \sigma_{\mathrm{post}}^2)$\;
  }
  Calculer $\mathrm{ELBO}_t$\;
}
\Retourner le r\'esum\'e post\'erieur
\caption{RSI : Inf\'erence d'\'Etat de R\`egles (VI champ-moyen)}
\label{alg:rsi}
\end{algorithm}

%% ─────────────────────────────────────────────────────────────────
\section{Garanties Th\'eoriques}
\label{sec:theory}

\subsection{T1 : Adaptabilit\'e r\'eglementaire en $O(1)$}

\begin{definition}[Mise \`a jour r\'eglementaire]
Une mise \`a jour $\mathcal{U}_k$ remplace $r_k$ par $r_k'$, diff\'erant
uniquement par les param\`etres $\Theta_k \to \Theta_k'$.
Le co\^ut $\mathcal{C}(\mathcal{U}_k)$ compte les op\'erations requises.
\end{definition}

\begin{theoreme}[Adaptabilit\'e $O(1)$]
\label{thm:o1}
Pour toute mise \`a jour $\mathcal{U}_k$, le co\^ut RSI est
$\mathcal{C}^{\mathrm{RSI}}(\mathcal{U}_k) = O(1)$, ind\'ependamment
de $|\mathbf{D}|$ et de $n$.
\end{theoreme}

\begin{proof}
L'a priori RSI se factorise en $P(\mathbf{S}) = \prod_i P(s_i | \Theta_i)$.
Apr\`es $\mathcal{U}_k$, le posterior mis \`a jour satisfait~:
\[
  P'(\mathbf{S} \mid \mathbf{D}) \;\propto\; P(\mathbf{S} \mid \mathbf{D})
  \cdot \underbrace{\frac{P(s_k' \mid \Theta_k')}{P(s_k \mid \Theta_k)}}_{%
    \text{scalaire, } O(1)}
\]
Ce ratio de correction ne d\'epend que des param\`etres de la r\`egle $k$.
Son calcul n\'ecessite l'\'evaluation de deux distributions
param\'etriques, soit $O(1)$ op\'erations.

\noindent\textbf{Stabilit\'e num\'erique.} Le ratio est appliqu\'e sur
la \emph{forme fonctionnelle} (le noyau) du posterior, et non sur sa
valeur absolue. La constante de normalisation est mise \`a jour
analytiquement gr\^ace aux familles conjugu\'ees (Beta-Bernoulli,
Gaussienne-Gaussienne), garantissant que le posterior reste une
distribution proprement normalis\'ee. Il n'y a aucun risque d'explosion
de variance ni de divergence num\'erique.
\end{proof}

\begin{corollaire}
Sur $T$ mises \`a jour r\'eglementaires~:
RSI~: $O(T)$~;
ML supervis\'e (r\'e-entra\^inement complet)~: $O(T \cdot |\mathbf{D}| \cdot E)$~;
PSL/MLN~: $O(T \cdot |\mathbf{D}| \cdot n)$.
RSI domine dans les environnements \`a forte fr\'equence r\'eglementaire.
\end{corollaire}

\subsection{T2 : Consistance de Bernstein-von Mises}

\begin{theoreme}[Consistance BvM]
\label{thm:bvm}
Soit $\mathbf{S}^*$ le vrai \'etat des r\`egles. Sous les conditions
(C1)~identifiabilit\'e, (C2)~$P(\mathbf{S}^*) > 0$,
(C3)~log-vraisemblance deux fois diff\'erentiable,
(C4)~information de Fisher $\mathcal{I}(\mathbf{S}^*)$ finie d\'efinie positive~:
\[
  \sqrt{m}\bigl(\hat{\mathbf{S}}_m - \mathbf{S}^*\bigr)
  \;\xrightarrow{d}\;
  \mathcal{N}\!\left(0,\;\mathcal{I}(\mathbf{S}^*)^{-1}\right)
  \quad (m \to \infty)
\]
\end{theoreme}

\begin{proof}[Preuve (esquisse)]
(i)~La log-vraisemblance normalis\'ee converge p.s. par la LGN.
(ii)~Un d\'eveloppement de Taylor au second ordre autour de $\mathbf{S}^*$
donne une approximation gaussienne.
(iii)~Sous (C2) et (C4), l'a priori est domin\'e par la vraisemblance
quand $m \to \infty$, ce qui donne la limite gaussienne
\citep{van2000asymptotic}.
\end{proof}

\begin{corollaire}[Robustesse \`a l'a priori]
Pour deux a priori RSI quelconques $P_1$, $P_2$~:
$\operatorname{TV}(P_1(\mathbf{S}|\mathbf{D}_m),\, P_2(\mathbf{S}|\mathbf{D}_m))
\to 0$ quand $m \to \infty$.
Deux experts avec des croyances initiales diff\'erentes convergent
vers la m\^eme \'evaluation de conformit\'e.
\end{corollaire}

\subsection{T3 : Convergence monotone de l'ELBO}

\begin{theoreme}[ELBO Monotone]
\label{thm:elbo}
Les mises \`a jour par ascension de coordonn\'ees RSI produisent une
s\'equence ELBO monotone non d\'ecroissante~:
$\mathrm{ELBO}_{t+1} \geq \mathrm{ELBO}_t$ pour tout $t \geq 0$.
\end{theoreme}

\begin{proof}
Chaque mise \`a jour minimise $\mathrm{KL}[Q_{\phi_i} \| P(s_i |
\mathbf{D})]$ en maintenant les autres facteurs fix\'es. Puisque
$\mathrm{ELBO} = -\mathrm{KL}[Q_\phi \| P(\mathbf{S}|\mathbf{D})] +
\log P(\mathbf{D})$ et que $\log P(\mathbf{D})$ est constant, chaque
mise \`a jour ne peut pas diminuer l'ELBO \citep{blei2017variational}.
\end{proof}

%% ─────────────────────────────────────────────────────────────────
\section{\'Evaluation Exp\'erimentale}
\label{sec:experiments}

\subsection{Le jeu de donn\'ees RSI-Togo-Fiscal-Synthetic}

Nous introduisons \textbf{RSI-Togo-Fiscal-Synthetic v1.0}, un benchmark
de 2~000 entreprises synth\'etiques ancr\'e dans les r\`egles officielles
de l'OTR~\citep{otr2024lf}, structur\'e en quatre couches~:

\begin{description}[leftmargin=1.5em]
  \item[Caract\'eristiques des entreprises] Secteur, r\'egion, CA d\'eclar\'e.
  \item[V\'erit\'e terrain latente] $(a_i, c_i, \delta_i)$ par r\`egle---
        l'\'etat non observ\'e que RSI doit inf\'erer.
  \item[Observations bruit\'ees] Taxes d\'eclar\'ees, proxys comportementaux
        (bancarisation, facturation \'electronique, d\'elais de paiement),
        avec ratio de sous-d\'eclaration moyen de 0,70 et 18--20\%
        de donn\'ees manquantes.
  \item[\'Etiquettes binaires] Labels de conformit\'e d\'eriv\'es pour
        l'\'evaluation des baselines.
\end{description}

\noindent Le dataset encode cinq r\`egles fiscales et un
\textbf{\'ev\'enement de changement r\'eglementaire document\'e}~:
le rel\`evement du seuil TVA de 60M \`a 100M FCFA (Loi n°2024-007
du 30 d\'ecembre 2024), permettant la validation empirique directe
du Th\'eor\`eme~\ref{thm:o1}.

\noindent\textbf{R\'ealisme du dataset.} La distribution des chiffres
d'affaires suit une loi log-normale par segment (TPU, RSI, r\'eel),
dont l'agr\'egation approxime une \emph{loi de puissance} (Power Law)
--- la signature universelle des distributions de revenus dans les
\'economies r\'eelles~\citep{gabaix2009power}. Les 2~000 entreprises
synth\'etiques ne sont pas uniformes~: elles reproduisent la structure
de concentration typique d'un march\'e africain (60\% d'informel,
25\% de PME, 15\% de grands contribuables), garantissant que les
r\'esultats ne sont pas des artefacts d'un dataset simpliste.

\subsection{Baselines et protocole exp\'erimental}

\begin{description}[leftmargin=1.5em]
  \item[SBR] Syst\`eme \`a Base de R\`egles d\'eterministe. Aucune
        incertitude. Fragile aux donn\'ees manquantes.
  \item[XGBoost] \citep{chen2016xgboost} Enti\`erement supervis\'e
        (150 estimateurs, profondeur 4).
  \item[MLP] Enti\`erement supervis\'e (64-32-16, ReLU, early stopping).
\end{description}

\noindent\textbf{Distinction critique.} RSI op\`ere en mode
\emph{z\'ero-shot}~: aucun exemple \'etiquet\'e de conformit\'e n'est
utilis\'e. XGBoost et MLP sont entra\^in\'es sous supervision compl\`ete.

\subsection{R\'esultats}

\paragraph{EXP-1 : Performance globale.}

\begin{table}[H]
\centering
\caption{Comparaison des performances. RSI n'utilise \emph{aucune \'etiquette}.
         $\dagger$~: supervision compl\`ete requise.
         \textbf{Gras}~: meilleure performance z\'ero-shot.}
\label{tab:main}
\renewcommand{\arraystretch}{1.2}
\begin{tabular}{lccccr}
\toprule
\textbf{Mod\`ele} & \textbf{F1} & \textbf{AUC} & \textbf{Rappel}
  & \textbf{\'Etiquettes ?} & \textbf{Co\^ut MAJ} \\
\midrule
\textbf{RSI (le n\^otre)} & \textbf{0,519} & \textbf{0,599} & \textbf{0,909}
  & \textbf{Non} & 0,2\,ms \\
SBR & 0,182 & --- & 0,214 & Non & $\sim$0\,ms \\
XGBoost$^\dagger$ & 0,967 & 0,999 & 0,978 & Oui & 407\,ms \\
MLP$^\dagger$ & 0,087 & 0,747 & 0,153 & Oui & 61\,ms \\
\bottomrule
\end{tabular}
\end{table}

RSI atteint F1\,=\,0,519 et AUC\,=\,0,599 sans aucune donn\'ee
d'entra\^inement \'etiquet\'ee, surpassant largement le syst\`eme
\`a base de r\`egles (F1\,=\,0,182). Le rappel \'elev\'e (0,909) est coh\'erent avec les besoins des
syst\`emes de surveillance fiscale, o\`u manquer un fraudeur
(\emph{faux n\'egatif}) est g\'en\'eralement plus co\^uteux
institutionnellement qu'une fausse alarme (\emph{faux positif}).
RSI s'adapte naturellement \`a cette asym\'etrie via le calibrage
du seuil de d\'ecision $\tau$, sans modifier l'algorithme d'inf\'erence.
XGBoost atteint un F1 quasi-parfait sous supervision compl\`ete, mais
ne peut pas op\'erer dans le cadre z\'ero-\'etiquette.

\paragraph{EXP-2 : Adaptabilit\'e $O(1)$ (Th\'eor\`eme~\ref{thm:o1}).}

\begin{table}[H]
\centering
\caption{Temps de mise \`a jour et F1 post-MAJ (p\'eriode 2025).}
\label{tab:update}
\renewcommand{\arraystretch}{1.2}
\begin{tabular}{lccc}
\toprule
\textbf{Mod\`ele} & \textbf{Temps MAJ} & \textbf{F1 post-MAJ}
  & \textbf{R\'e-entra\^inement ?} \\
\midrule
RSI (le n\^otre) & \textbf{0,2\,ms} & 0,411 & Non \\
XGBoost & 407\,ms & 0,954 & Oui \\
MLP & 61\,ms & 0,000 & Oui \\
\bottomrule
\end{tabular}
\end{table}

RSI se met \`a jour en 0,2\,ms---une acc\'el\'eration de
$\mathbf{2~000\times}$ par rapport \`a XGBoost. Sur $T$ mises \`a jour
annuelles (Togo~: $T \approx 4$--$6$), les \'economies cumulatives
croissent lin\'eairement en $T$ et exponentiellement avec la taille
du dataset, confirmant le Corollaire~1.

\paragraph{EXP-3 : Consistance BvM (Th\'eor\`eme~\ref{thm:bvm}).}
L'incertitude post\'erieure RSI $\sigma[c_i | \mathbf{D}]$ d\'ecro\^it
monotonement avec la taille d'\'echantillon $N$, conforme au
Th\'eor\`eme~\ref{thm:bvm}. Pour $N < 25$, RSI atteint un rappel
comp\'etitif tandis que XGBoost ne peut pas fonctionner (exemples
\'etiquet\'es insuffisants).

\paragraph{EXP-4 : Robustesse aux donn\'ees manquantes.}
Sous 20\% de donn\'ees manquantes (typique pour l'Afrique francophone),
le F1 de RSI se d\'egrade de seulement $-0,003$ ($< 1\%$), tandis que
le SBR perd $> 22\%$ de F1. \`A 50\% de donn\'ees manquantes, RSI
reste fonctionnel (F1\,=\,0,521) alors que les m\'ethodes
d\'eterministes s'effondrent.

\paragraph{EXP-5 : Convergence ELBO (Th\'eor\`eme~\ref{thm:elbo}).}
La s\'equence ELBO est empiriquement monotone non d\'ecroissante et
converge en 7 it\'erations, confirmant le Th\'eor\`eme~\ref{thm:elbo}.

\subsection{Interpr\'etabilit\'e : un exemple concret}

Pour une entreprise avec un CA d\'eclar\'e de 72M FCFA et aucune
TVA d\'eclar\'ee, RSI produit~:
\begin{itemize}[leftmargin=2em]
  \item $P(\text{TVA active} \mid D) = 0,91$ --- la r\`egle s'applique
        tr\`es probablement
  \item $\mathbb{E}[c_{\text{TVA}} \mid D] = 0,23 \pm 0,18$ --- faible
        conformit\'e, forte incertitude
  \item $\mathbb{E}[\delta_{\text{TVA}} \mid D] = +2,4$M FCFA --- d\'erive
        positive du seuil
\end{itemize}
Un auditeur agit directement sur cet output.
\emph{Le posterior est l'explication}---aucun outil post-hoc (SHAP,
LIME) n'est n\'ecessaire.

%% ─────────────────────────────────────────────────────────────────
\section{Discussion}
\label{sec:discussion}

\paragraph{G\'en\'eralisabilit\'e.}
RSI est agnostique au domaine. Les extensions imm\'ediates incluent
la conformit\'e aux protocoles m\'edicaux, la lutte anti-blanchiment,
la surveillance environnementale et le suivi des contrats juridiques.
Tout domaine satisfaisant la D\'efinition~1 est candidat.

\paragraph{Limites.}
La qualit\'e de l'inf\'erence d\'epend de la pr\'ecision des
hyperparams de l'a priori. Des a priori mal sp\'ecifi\'es biaisent
les r\'esultats, bien que T2 garantisse la r\'ecup\'eration asymptotique.
Les fonctions de vraisemblance segment\'ees n\'ecessitent actuellement
une expertise du domaine~; des travaux futurs exploreront l'inf\'erence
amortie pour l'apprentissage de la vraisemblance.

\paragraph{Relation avec le filtrage de Kalman.}
RSI partage des parall\`eles architecturaux avec le filtre de Kalman~:
l'a priori joue le r\^ole du mod\`ele de processus, la vraisemblance
est le mod\`ele de mesure, et le posterior est l'estimation d'\'etat.
Cette connexion ouvre des voies pour RSI s\'equentiel, suivant
l'\'evolution des \'etats de r\`egles dans le temps.

\paragraph{Passage \`a l'\'echelle.}
La VI champ-moyen s'\'echelonne lin\'eairement en $n$. Pour les codes
fiscaux complets (centaines de r\`egles), des extensions en graphes
de facteurs via la propagation de croyances g\`erent les d\'ependances
entre r\`egles.

%% ─────────────────────────────────────────────────────────────────
\section{Conclusion}
\label{sec:conclusion}

Nous avons introduit \textbf{RSI (Rule-State Inference)}, un cadre
bay\'esien qui inverse formellement le paradigme dominant de la
surveillance de la conformit\'e. En traitant les r\`egles r\'eglementaires
autorit\'ees comme des a priori structur\'es et en formulant l'\'evaluation
de conformit\'e comme une inf\'erence post\'erieure sur un espace
d'\'etats latents, RSI offre ce qu'aucun cadre existant ne fournit~:
surveillance de conformit\'e z\'ero-shot, adaptabilit\'e r\'eglementaire
$O(1)$, quantification native de l'incertitude, et pleine
interpr\'etabilit\'e.

Trois th\'eor\`emes \'etablissent les fondements de RSI~: adaptabilit\'e
$O(1)$ (T1), consistance de Bernstein-von Mises (T2), et convergence
monotone de l'ELBO (T3). Les exp\'eriences sur RSI-Togo-Fiscal-Synthetic
v1.0---ancr\'e dans la loi fiscale togolaise r\'eelle---valident les
trois empiriquement et d\'emontrent une acc\'el\'eration de $2~000\times$
par rapport aux m\'ethodes bas\'ees sur le r\'e-entra\^inement.

RSI est con\c{c}u pour les r\'ealit\'es des environnements r\'egis
par des r\`egles dans les pays du Sud~: changements r\'eglementaires
fr\'equents, raret\'e structurelle des donn\'ees, et exigences
institutionnelles de d\'ecisions auditable et interpr\'etables.
Nous publions le cadre complet, le dataset et l'impl\'ementation
pour soutenir la recherche reproductible.

\medskip
\noindent\textbf{Travaux futurs~:} RSI s\'equentiel pour le suivi
longitudinal de la conformit\'e, extensions multi-pays pour la zone
CEDEAO, et vraisemblances apprises par inf\'erence variationnelle amortie.

%% ─────────────────────────────────────────────────────────────────
\section*{Remerciements}

L'auteur exprime sa profonde gratitude au \textbf{Togo DataLab}
(Minist\`ere de l'\'Economie Num\'erique et de la Transformation
Digitale) ainsi qu'\`a l'\textbf{Office Togolais des Recettes (OTR)}
pour leur accueil, l'acc\`es aux probl\'ematiques m\'etier et leur
expertise in\'estimable sur le cadre r\'eglementaire fiscal togolais.

Ce travail a \'egalement b\'en\'efici\'e de l'environnement acad\'emique
de l'\textbf{Institut National des Postes et T\'el\'ecommunications
(INPT)} de Rabat. L'auteur tient \`a remercier ses collaborateurs et
mentors pour les discussions enrichissantes qui ont permis de consolider
les bases th\'eoriques de ce cadre.

\textit{Note~: Les opinions exprim\'ees dans ce document sont celles
de l'auteur et ne refl\`etent pas n\'ecessairement les positions
officielles des institutions mentionn\'ees.}

\medskip
\noindent\textit{Déclaration IA : L'auteur a utilisé des outils 
assistés par IA pour la rédaction du manuscrit et la mise en forme 
\LaTeX{}. Toutes les contributions scientifiques, affirmations 
théoriques, conception expérimentale et conclusions relèvent de 
la seule responsabilité de l'auteur.}

%% ─────────────────────────────────────────────────────────────────
\bibliographystyle{plainnat}
\begin{thebibliography}{99}

\bibitem[Bach et al.(2017)]{bach2017hinge}
Bach, S., Broecheler, M., Huang, B., et Getoor, L. (2017).
Hinge-loss Markov random fields and probabilistic soft logic.
\textit{Journal of Machine Learning Research}, 18(109):1--67.

\bibitem[Badreddine et al.(2022)]{badreddine2022logic}
Badreddine, S., d'Avila Garcez, A., Serafini, L., et Spranger, M. (2022).
Logic tensor networks.
\textit{Artificial Intelligence}, 303:103649.

\bibitem[Blei et al.(2017)]{blei2017variational}
Blei, D. M., Kucukelbir, A., et McAuliffe, J. D. (2017).
Variational inference: A review for statisticians.
\textit{Journal of the American Statistical Association}, 112(518):859--877.

\bibitem[Chen et Guestrin(2016)]{chen2016xgboost}
Chen, T. et Guestrin, C. (2016).
{XGBoost} : A scalable tree boosting system.
In \textit{Actes de KDD}, pp.\ 785--794.

\bibitem[Fr{\'e}nay et Verleysen(2014)]{frenay2014classification}
Fr{\'e}nay, B. et Verleysen, M. (2014).
Classification in the presence of label noise: a survey.
\textit{IEEE Transactions on Neural Networks and Learning Systems}, 25(5):845--869.

\bibitem[Garc{\'e}z et Lamb(2022)]{garcez2022neural}
Garc{\'e}z, A. d'Avila et Lamb, L. C. (2022).
Neurosymbolic AI: The third wave.
\textit{Artificial Intelligence Review}.

\bibitem[Gabaix(2009)]{gabaix2009power}
Gabaix, X. (2009).
Power laws in economics and finance.
\textit{Annual Review of Economics}, 1(1):255--294.

\bibitem[Gomes et al.(2022)]{gomes2022machine}
Gomes, M., Batista, J., et Carvalho, J. (2022).
Machine learning for tax gap estimation.
\textit{Expert Systems with Applications}, 198:116810.

\bibitem[Letham et al.(2015)]{letham2015interpretable}
Letham, B., Rudin, C., McCormick, T. H., et Madigan, D. (2015).
Interpretable classifiers using rules and Bayesian analysis.
\textit{Annals of Applied Statistics}, 9(3):1350--1371.

\bibitem[OTR(2024)]{otr2024lf}
Office Togolais des Recettes (2024).
Loi n°2024-007 du 30 d{\'e}cembre 2024 portant loi de finances
pour l'exercice 2025. Lom{\'e}, Togo.

\bibitem[Richardson et Domingos(2006)]{richardson2006markov}
Richardson, M. et Domingos, P. (2006).
Markov logic networks.
\textit{Machine Learning}, 62(1--2):107--136.

\bibitem[Rockt{\"a}schel et Riedel(2017)]{rocktaschel2017end}
Rockt{\"a}schel, T. et Riedel, S. (2017).
End-to-end differentiable proving.
In \textit{NeurIPS}, pp.\ 3788--3800.

\bibitem[van der Vaart(2000)]{van2000asymptotic}
van der Vaart, A. W. (2000).
\textit{Asymptotic Statistics}.
Cambridge University Press.

\end{thebibliography}

\end{document}